% nlt-worked-examples.tex — NLT 算例
% Source: examples/nlt-worked-examples.md

\chapter{NLT 算例}\label{ex:nlt}

对应章节：第 \ref{ch:nlt} 章（NLT）。

\begin{examplebox}[title={例 1：NLT None（线性变换）}]

\textbf{输入条件}：单个像素值 $x = 200$，输入位深 $d = 8$，$B_w = 20$。

\medskip
\textbf{已知参数}：

\begin{tabular}{@{}lc@{}}
\toprule
参数 & 值 \\
\midrule
$d$ & 8 \\
$B_w$ & 20 \\
$s = B_w - d$ & 12 \\
$\mathit{dclev} = 2^{B_w-1}$ & $2^{19} = 524288$ \\
$\mathit{dclev\_and\_rounding} = \mathit{dclev} + 2^{s-1}$ & $524288 + 2048 = 526336$ \\
$\mathit{max\_val} = 2^d - 1$ & 255 \\
\bottomrule
\end{tabular}

\medskip
\textbf{前向变换}（对照 \texttt{nlt\_forward\_linear()}）：

$v_{\mathit{out}} = (200 \ll 12) - 524288 = 819200 - 524288 = 294912$

\textbf{逆向变换}（对照 \texttt{nlt\_inverse\_linear()}）：

$x_{\mathit{rec}} = \mathrm{clip}((294912 + 526336) \gg 12,\; 0,\; 255) = \mathrm{clip}(200.5,\; 0,\; 255) = 200$

\medskip
\textbf{中间变量表}：

\begin{tabular}{@{}cllll@{}}
\toprule
步骤 & 方向 & 操作 & 计算 & 结果 \\
\midrule
1 & 前向 & 左移 $s$ 位 & $200 \times 4096$ & 819200 \\
2 & 前向 & 减 $\mathit{dclev}$ & $819200 - 524288$ & 294912 \\
3 & 逆向 & 加 $\mathit{dclev}$ + 舍入 & $294912 + 526336$ & 821248 \\
4 & 逆向 & 右移 $s$ 位 & $821248 / 4096$ & $200.5 \to 200$ \\
\bottomrule
\end{tabular}

\medskip
\textbf{最终结果}：输入 200，经前向 $\to$ 逆向变换后恢复为 200。NLT 线性变换在整数域内无损。
\end{examplebox}

\begin{examplebox}[title={例 2：NLT Quadratic（二次变换）}]

\textbf{输入条件}：单个像素值 $x = 100$，输入位深 $d = 8$，$B_w = 10$，参数 $\sigma = 0, \alpha = 0$。

\medskip
\textbf{已知参数}：

\begin{tabular}{@{}lc@{}}
\toprule
参数 & 值 \\
\midrule
$d$ & 8 \\
$B_w$ & 10 \\
$s = B_w - d$ & 2 \\
$\mathit{dclev} = 2^{B_w-1}$ & 512 \\
$\mathit{max\_coef} = 2^{B_w} - 1$ & 1023 \\
$v_{\mathit{dco}} = \alpha - \sigma \times 32768$ & 0 \\
\bottomrule
\end{tabular}

\medskip
\textbf{前向变换}（对照 \texttt{nlt\_forward\_quadratic()}）：

\begin{enumerate}
  \item $v_1 = \mathrm{clip}(100 - 0,\; 0,\; 255) = 100$
  \item $v_2 = 100 \ll 2 = 400$
  \item $\text{sqrt\_approx\_fixpoint}(400, 10) \to r = 40$，$r \gg 1 = 20$
  \item $v_{\mathit{out}} = 20 - 512 = -492$
\end{enumerate}

\textbf{逆向变换}（对照 \texttt{nlt\_inverse\_quadratic()}）：

\begin{enumerate}
  \item $v = \mathrm{clip}(-492 + 512,\; 0,\; 1023) = 20$
  \item $v^2 = 20^2 = 400$
  \item $s_{\mathit{inv}} = 2 \cdot B_w - d = 12$
  \item $x_{\mathit{rec}} = \mathrm{clip}((400 + (1 \ll 11)) \gg 12,\; 0,\; 255) = 0$
\end{enumerate}

\begin{engineeringnote}
逆变换中 $s_{\mathit{inv}} = 12$，$v^2 = 400$ 远小于 $2^{12} = 4096$，右移后结果为 0。这说明 Quadratic NLT 的正逆变换不是严格互逆的——前向使用定点 $\sqrt{\cdot}$，逆向使用整数平方后右移，小值会因精度不足丢失信息。这是 JPEG XS 作为有损压缩标准的设计选择。
\end{engineeringnote}

\textbf{最终结果}：$B_w = 10$ 时，前向 $v_{\mathit{out}} = -492$，逆向 $x_{\mathit{rec}} = 0$。定点平方根本身精确（$\sqrt{400} = 20$），但逆变换的整数右移导致小值丢失。
\end{examplebox}

\begin{examplebox}[title={例 3：NLT Extended（扩展三段式变换）}]

\textbf{输入条件}：$B_w = 10$，$E = 3$，$d = 8$，$T_1 = 84$，$T_2 = 256$。

\medskip
\textbf{派生常量}（对照 \texttt{nlt\_forward\_extended()}）：

\begin{tabular}{@{}clll@{}}
\toprule
常量 & 公式 & 计算 & 结果 \\
\midrule
$B_2$ & $T_1^2$ & $84^2$ & 7056 \\
$A_1$ & $B_2 + T_1 \times 2^e + 2^{2e-2}$ & $7056 + 10752 + 4096$ & 21904 \\
$B_1$ & $T_1 + 2^{e-1}$ & $84 + 64$ & 148 \\
$A_3$ & $B_2 + T_2 \times 2^e - 2^{2e-2}$ & $7056 + 32768 - 4096$ & 35728 \\
$B_3$ & $T_2 - 2^{e-1}$ & $256 - 64$ & 192 \\
$Q_1$ & $B_2 + T_1 \times 2^e$ & $7056 + 10752$ & 17808 \\
$Q_2$ & $B_2 + T_2 \times 2^e$ & $7056 + 32768$ & 39824 \\
\bottomrule
\end{tabular}

\medskip
\textbf{子例 3a：线性区}（$x = 8$）：

$v = 8 \ll 12 = 32768$。$Q_1 = 17808 \leq 32768 < 39824 = Q_2$，落入线性区。

前向：$v_{\mathit{out}} = (32768 - 7056) / 128 - 512 = 201 - 512 = -311$

逆向：$v = -311 + 512 = 201$，$201 \times 128 + 7056 = 32784$，$32784 \gg 12 = 8$。

\textbf{往返验证：$x = 8 \to v_{\mathit{out}} = -311 \to x_{\mathit{rec}} = 8$。无损。}

\medskip
\textbf{子例 3b：高区}（$x = 100$）：

$v = 100 \ll 12 = 409600 \geq Q_2$，落入高区。

前向：$v = B_3 + \sqrt{409600 - 35728} = 192 + 611 = 803$，$v_{\mathit{out}} = 803 - 512 = 291$

逆向：$v - B_3 = 611$，$611^2 = 373321$，$(35728 + 373321) \gg 12 = 99$

\textbf{往返结果：$x = 100 \to v_{\mathit{out}} = 291 \to x_{\mathit{rec}} = 99$。偏差 1，由 \texttt{sqrt()} 浮点截断引入。}

\medskip
\textbf{子例 3c：低区}（$x = 3$）：

$v = 3 \ll 12 = 12288 < Q_1$，落入低区。

前向：$v = B_1 - \sqrt{21904 - 12288} = 148 - 98 = 50$，$v_{\mathit{out}} = 50 - 512 = -462$

逆向：$B_1 - v = 98$，$98^2 = 9604$，$(21904 - 9604) \gg 12 = 2$

\textbf{往返结果：$x = 3 \to v_{\mathit{out}} = -462 \to x_{\mathit{rec}} = 2$。偏差 1。}

\medskip
\textbf{最终结果}：

\begin{tabular}{@{}clccl@{}}
\toprule
像素值 $x$ & 区域 & $v_{\mathit{out}}$ & $x_{\mathit{rec}}$ & 偏差 \\
\midrule
8 & 线性区 & $-311$ & 8 & 0（无损） \\
3 & 低区 & $-462$ & 2 & 1 \\
100 & 高区 & 291 & 99 & 1 \\
\bottomrule
\end{tabular}

线性区往返无损。低区和高区各引入 1 个灰度级偏差，来源于 \texttt{sqrt()} 浮点截断。
\end{examplebox}

\begin{examplebox}[title={例 4：Extended NLT 自动阈值生成}]

\textbf{输入条件}：目标阈值 $\mathit{th}_1 = 10$，$\mathit{th}_2 = 100$（像素域），$B_w = 18$，bpp $= 8$。

\medskip
\textbf{计算 $T_1$}（对照 \texttt{xs\_config\_nlt\_extended\_auto\_thresholds()}）：

\begin{align*}
T_1 &= \sqrt{2^{2e-2} + \mathit{th}_1 \times 2^s} - 2^{e-1} \\
&= \sqrt{268435456 + 10 \times 268435456} - 16384 \\
&= \sqrt{2952790016} - 16384 = 54339 - 16384 = 37955
\end{align*}

\textbf{计算 $T_2$}：

\begin{align*}
T_2 &= (\mathit{th}_2 \times 2^s - T_1^2) \gg e \\
&= (100 \times 268435456 - 1440582025) \gg 15 \\
&= 25402963575 \gg 15 = 775375
\end{align*}

\textbf{区域分布}：

\begin{tabular}{@{}clc@{}}
\toprule
区域 & 像素范围 & 说明 \\
\midrule
低区 & 0--9 & 暗部，二次压缩 \\
线性区 & 10--99 & 中间调，线性映射 \\
高区 & 100--255 & 亮部，二次扩展 \\
\bottomrule
\end{tabular}

给定 $\mathit{th}_1 = 10, \mathit{th}_2 = 100$，自动生成 $T_1 = 37955$、$T_2 = 775375$。
\end{examplebox}
